% Generated draft: requires mathematical and source review.
\documentclass[11pt]{article}
\usepackage[margin=25mm]{geometry}
\usepackage{fontspec}
\setmainfont{Latin Modern Roman}
\usepackage{amsmath,amssymb,mathtools,unicode-math}
\setmathfont{Latin Modern Math}
\usepackage{enumitem,needspace,xurl,hyperref}
\setlength{\parindent}{0pt}
\setlength{\parskip}{5pt}
\setlength{\emergencystretch}{3em}
\newcommand{\N}{\mathbb{N}}
\newcommand{\Z}{\mathbb{Z}}
\newcommand{\Q}{\mathbb{Q}}
\newcommand{\R}{\mathbb{R}}
\newcommand{\C}{\mathbb{C}}
\newcommand{\F}{\mathbb{F}}
\newcommand{\A}{\mathbb{A}}
\newcommand{\PP}{\mathbb{P}}
\newcommand{\E}{\mathbb{E}}
\newcommand{\eps}{\varepsilon}
\newcommand{\dd}{\,\mathrm{d}}
\begin{document}
\section*{UnsolvedMath ID: 3000027}
\textbf{Source problem code:} catalogue entry
\subsection{Definitions and mathematical statement}
Let $G=(V,E)$ be an undirected graph with root $r\in V$. For a graph $H$ and integer $q>0$, $q$-edge-connectivity means
\[
|\delta_H(A)|\geq q
\quad\text{for every nonempty proper }A\subset V(H),
\]
where $\delta_H(A)$ is the edge cut. Retaining the source's definition, $G$ is rooted $(k,l)$-connected if, for every $X\subseteq V\setminus\{r\}$,
\[
G-X\text{ is }(k-|X|)l\text{-edge-connected}.
\]
In the intended positive-integer parameter range, nonpositive thresholds are naturally vacuous; equivalently only $|X|<k$ imposes a positive cut requirement. That convention is not stated in the source and must be checked.

The question is whether every such $G$ contains subgraphs
\[
G_i=(V,E_i)\quad(1\leq i\leq l),\qquad
E_i\subseteq E,\quad E_i\cap E_j=\varnothing\ (i\ne j),
\]
with each $G_i$ spanning and rooted $k$-connected at the same root $r$. ``Contains'' requests disjoint edge sets, without explicitly requiring their union to exhaust $E$.

The source does not define rooted $k$-connected. A usual root-vertex-connectivity reading is that for each $v\ne r$ there are $k$ paths from $r$ to $v$ whose internal vertices are pairwise disjoint. This is distinct from $k$ edge-disjoint paths and must be verified as the intended target. The source also leaves finite versus infinite graphs, parallel-edge conventions, and small-order/parameter boundary cases unstated. No such choice is silently used to turn the conjecture into a known edge-connectivity decomposition theorem. The supplied summary places the problem in root-connectivity/branching decomposition theory and reports no verified resolution.
\subsection{Short English statement}
The conjecture asks whether the stated edge-connectivity conditions after deleting nonroot vertices force $l$ disjoint spanning subgraphs, each $k$-connected from the root in the intended sense. The deletion condition is explicit, but the target root-connectivity and graph conventions require source review.
\subsection{Sources}
\begin{itemize}
\item Local UnsolvedMath registry, permanent ID 3000027, exact source code catalogue entry, ``Decomposing rooted $(k,l)$-connected graphs into rooted $k$-connected parts''; status open, classification OPEN-TRIAGE.
\item Egres Open, \emph{List of open problems}, indexed page ``Decomposing rooted (k,l)-connected graphs into rooted k-connected parts,'' \url{https://oldlemon.cs.elte.hu/egres/open/List_of_open_problems}, HTML extraction accessed July 29, 2026. The record says the archived index classes the page as open but is not a current-resolution authority. No separate problem-page URL or primary paper is supplied.
\item The supplied summary reports openness as far as verifiable there and no verified resolution in the branching/arborescence decomposition literature, without a precise theorem or further citation.
\end{itemize}
\end{document}
